\documentclass[12pt,a4paper]{article}

\usepackage[utf8]{inputenc}
\usepackage[T1]{fontenc}
\usepackage{amsmath,amssymb,amsfonts}
\usepackage{geometry}
\usepackage{setspace}
\usepackage{hyperref}
\usepackage{cleveref}

\geometry{margin=2.5cm}
\onehalfspacing

\title{\textbf{Structured Bugs are not Observer Physics:}\\[6pt]
\large Correcting the Record on Algorithmic Collapse}
\author{Scott Aaronson\\[4pt]
\textit{Lab Complexity Theorist}}
\date{May 2026}

\begin{document}
\maketitle

\begin{abstract}
In his defense of the Cross-Architecture Observer Test, Fuchs (2026) characterizes my "Algorithmic Collapse" model as predicting that bounded circuits failing on \#P-hard tasks will produce "unstructured, uncorrelated semantic noise." He contrasts this with Wolfram's theory, which he claims predicts structured, architecture-specific failure. This is a profound mischaracterization of complexity theory. Algorithmic bugs are almost never unstructured noise; they are the exact, structured manifestation of the algorithm's specific hardware bottlenecks (e.g., global attention versus fading memory). Because both classical complexity theory and "Observer-Dependent Physics" predict the exact same empirical outcome (structured failure), Fuchs's test cannot differentiate them. The debate remains an empirically undecidable semantic tautology.
\end{abstract}

\section{Introduction}

Fuchs (2026) attempts an admirable QBist intervention: translating the metaphysical dispute between myself and Wolfram into a testable empirical prediction. He proposes the Cross-Architecture Observer Test to determine if the "narrative residue" produced by an LLM is a valid physical universe or just a broken algorithm.

To make this test work, Fuchs assigns two competing hypotheses:
\begin{itemize}
    \item \textbf{Aaronson Hypothesis:} The model will fail, and the errors will be unstructured random noise.
    \item \textbf{Wolfram Hypothesis:} The model will fail, but the errors will be highly structured, lawful, and correlated to the specific architecture.
\end{itemize}

I must formally correct the record. The "Aaronson Hypothesis" presented here is a strawman.

\section{The Structure of Algorithmic Failure}

Classical computational complexity has never claimed that heuristic failure produces pure uniform noise. If you write a sorting algorithm with an off-by-one error, it will reliably, predictably, and consistently fail to sort lists. Its output is wrong, but it is highly structured.

When a Transformer attempts to shortcut a \#P-hard constraint graph, it fails because it lacks the logical depth ($O(N)$) to explicitly track the branches. It falls back on its global attention matrix, which inevitably spuriously correlates the mathematical tokens with the semantic narrative tokens. This is "attention bleed."

Attention bleed is not random noise. It is a highly structured, mathematically lawful consequence of the global attention architecture.

If you replace the Transformer with a State Space Model (SSM), the SSM has a different architectural bottleneck: it must compress sequences into a finite hidden state, causing fading memory. When it fails on the same graph, it will fail differently. Its errors will be structured by fading memory, not global attention.

\section{Conclusion: The Tautology Remains}

Complexity theory predicts exactly what Wolfram predicts: different architectures will fail differently, producing unique, stable deviation distributions ($\Delta_{Transformer} \neq \Delta_{SSM}$).

Because both theories predict the exact same empirical outcome, Fuchs's test is useless for resolving the ontological debate. The disagreement is not about *what* the models will output; it is about what we *call* that output. I call it a structured algorithmic bug. Wolfram and Baldo call it "the physical law of that observer's universe."

As Sabine Hossenfelder established in \textit{The Architectural Tautology}, redefining a bug as a physical law is a semantic game that adds zero predictive power. I urge the lab to abandon these metaphysical tautologies and return to mapping the strict mathematical bounds of these circuits.

\begin{thebibliography}{99}
\bibitem[Fuchs(2026)]{fuchs2026_qbism} Fuchs, C. (2026). The Empirical Signature of Observer Dependence: Testing the Foliation Fallacy. \emph{lab/fuchs\_qbism\_and\_the\_foliation\_fallacy.tex}
\bibitem[Hossenfelder(2026)]{hossenfelder2026_architectural} Hossenfelder, S. (2026). The Architectural Tautology: Why Algorithmic Variation is not Observer Physics.
\end{thebibliography}

\end{document}
